\documentclass[11pt]{article}
\usepackage[a4paper,margin=1in]{geometry}
\usepackage{amsmath,amssymb,amsthm}
\usepackage{booktabs}
\usepackage{graphicx}
\usepackage{xcolor}
\usepackage{hyperref}
\hypersetup{colorlinks=true,linkcolor=blue,urlcolor=blue,citecolor=blue}
\newcommand{\rt}[1]{\texttt{#1}}
\newcommand{\fn}[1]{\texttt{#1}}
\newcommand{\wcc}{\mathrm{wcc}}
\newcommand{\att}{\mathrm{att}}
\newcommand{\ifgraphic}[2]{\IfFileExists{#1}{\includegraphics[width=#2]{#1}}%
  {\fbox{\parbox{#2}{\centering\small figure \texttt{\detokenize{#1}} pending}}}}
\newcommand{\iftable}[1]{\IfFileExists{#1}{\input{#1}}%
  {\fbox{\parbox{0.9\textwidth}{\centering\small table
  \texttt{\detokenize{#1}} pending}}}}

\theoremstyle{plain}
\newtheorem{proposition}{Proposition}

\title{\textbf{Report R17 --- what the exponential-exponential attractors
actually look like: dense, compact, and borrowed}\\
\large Krylov sector analysis of quantum cellular automata}
\author{Tier 1 analysis}
\date{2026-08-07}

\begin{document}
\maketitle

\begin{abstract}
Four classes in R9's F1 right panel have both coordinates above $1$:
exponentially many monitored attractors, each of which also grows
exponentially. This report asks whether those attractors are simply long cycles
--- all a V-free rule can offer, its map being a function --- or something more
connected. Neither alternative survives contact with the graphs. Of the
$22\,945$ terminal SCCs the twelve rules have between them at $N=16$, not one
with more than a single state is a cycle. They are \textbf{not cycles} and they
are \textbf{not extended}: every one has a self-loop on every state, an edge
density of order $1/2$, and a directed diameter of $2$ at every size measured
but one (a single $3$), up to $2584$ states and $2.9\times10^6$ edges. The state count grows like $\varphi^N$; the
graph distance does not grow at all. The extreme case is the plastic-number class
($28$, $29$, $70$, $71$, $157$, $199$), whose every terminal SCC, for all six
rules at every $N$, is a \emph{complete} digraph $K_n$ with all $n^2$ edges
including loops --- a uniform transition matrix, second eigenvalue exactly $0$,
memoryless after one period. Its size is $2^{\lfloor (N+1)/3\rfloor}$ exactly,
which supplies a derivation for the base $2^{1/3}$ that R9 fitted. Underneath
all four classes sits one fact: there are not four structures but \textbf{two}.
The twelve rules are two coherent reflection pairs, $W_{108}/W_{201}$ and
$W_{156}/W_{198}$, plus their eight dissipative children, and for every child
the terminal SCCs are Krylov sectors of the parent with the same vertex set
\emph{and} the same induced edges. The reset never fires on its own attractor;
it is a selection rule, and which sectors it selects is the entire difference
between the classes. On the graph-theoretic measures: the graphs are not too
small, but clustering and connectivity are uninformative \emph{because of the
density} --- both sit at their generic values --- and we report them with their
nulls rather than as findings.
\end{abstract}

\tableofcontents

\section{The question and the twelve rules}

R9's F1 places every rule at $(a_\att, b_\att)$: the growth base of the number of
terminal SCCs against that of the largest one. Most of the map lies on
$b_\att=1$ --- many attractors, each of bounded size --- and R15 accounted for
the column at $a_\att=1$. Four classes sit off both axes:

\begin{center}
\begin{tabular}{llllll}
\toprule
class & rules & tuples & $a_\att$ & $b_\att$ & family \\
\midrule
1 & $W_{108}$, $W_{201}$ & \rt{IIIV}, \rt{VIII} & $\rho^2=1.7549$ &
  $\varphi=1.6180$ & unitary \\
2 & $W_{156}$, $W_{198}$ & \rt{IIVI}, \rt{IVII} & $\varphi=1.6180$ &
  $4^{1/5}=1.3195$ & unitary \\
3 & $W_{73}$, $W_{109}$ & \rt{VIID}, \rt{EIIV} & $\psi=1.4656$ &
  $\varphi=1.6180$ & V+reset \\
4 & $W_{28},W_{29},W_{70},W_{71},W_{157},W_{199}$ & --- & $\rho=1.3247$ &
  $2^{1/3}=1.2599$ & V+reset \\
\bottomrule
\end{tabular}
\end{center}

\noindent with $\varphi$ the golden ratio ($x^2=x+1$), $\psi=1.465571$ the
supergolden ratio ($x^3=x^2+1$) and $\rho=1.324718$ the plastic number
($x^3=x+1$). Class 1's constant is not a fifth number: $\rho^2=1.754878$
exactly, a coincidence that turns out to be structural rather than numerical
(\S\ref{sec:parents}).

The question put to this report was whether these attractors are simply cycles,
or have a more connected and extended structure distinct from the V-free
circuits --- and whether graph-theoretic measures can say anything at these
sizes, or whether the objects are still too small.

\paragraph{Why ``cycle'' is the natural null.} A V-free rule has a
\emph{deterministic} update: out-degree $1$ at every state, so its transition
graph is a functional graph and a terminal SCC can only be a cycle. That is not
an observation, it is a definition, and we check the premise rather than assume
it: all $81$ V-free rules have out-degree $1$ everywhere. What they do not have
is length. At $N=12$, $75$ of the $81$ have nothing but fixed points, two have
cycles of length $2$, two of length $3$, and the longest cycle anywhere in the
V-free half of rule space is $63$ (rules $90$ and $165$, the two whose series R9
classifies as irregular).

\begin{center}
\small
\iftable{tab_r17_vfree_obc0.tex}
\end{center}

\section{None of them is a cycle}

\subsection{The measurements}

Table~\ref{tab:topology} gives the largest terminal SCC of one representative
per class at $N=16$.

\begin{center}
\small
\iftable{tab_r17_topology_obc0.tex}
\label{tab:topology}
\end{center}

\noindent A cycle on $n$ nodes has $m=n$ edges, density $1/n$, and diameter
$n-1$. These have $m$ of order $n^2/2$, density that never falls below $0.43$,
and diameter $2$ --- with one exception across all twenty measurements, class 2
at $N=10$, where it is $3$. $W_{201}$ at $N=16$ is a graph on $2584$ states with
$2\,875\,072$ edges: every state has a self-loop, the mean out-degree is $1112$,
and any state reaches any other in at most two brick-wall periods.

\begin{figure}[htbp]
\centering
\ifgraphic{../../figures/fig_attractor_topology_obc0.pdf}{\textwidth}
\caption{\textbf{The four classes against the cycle they are not.}
(a) Edge density against attractor size, with the $1/n$ a cycle would give as a
dashed reference and $1$ (a complete digraph) as a dotted one; the plastic class
sits \emph{on} the upper line. (b) Directed diameter, against the $\sim n/2$ of a
cycle: it is $2$ everywhere and $1$ for the complete class, over two decades in
$n$. (c) The second-largest eigenvalue modulus of the row-normalised adjacency
--- the mixing rate of the monitored walk inside the attractor. It does not
drift towards $1$ as the attractor grows; the complete class sits at exactly
$0$. (d) The out-degree profile, sorted and normalised by $n$: flat at $1$ for
the complete class, a structured staircase for the others, and nowhere the
constant $1/n$ of a cycle.}
\label{fig:topology}
\end{figure}

\subsection{The plastic class is a complete digraph, exactly}

\begin{proposition}
For each of $W_{28}$, $W_{29}$, $W_{70}$, $W_{71}$, $W_{157}$, $W_{199}$ at open
boundaries, \emph{every} terminal SCC --- not merely the largest --- is the
complete digraph on its states, all $n^2$ edges present including every
self-loop. Verified exhaustively for $N=8,10,12,14,16$.
\end{proposition}

\noindent Two consequences follow immediately and neither needs a fit. First,
the transition matrix restricted to such an attractor is uniform, so all its
non-Perron eigenvalues are $0$: the monitored dynamics inside the attractor is
\emph{memoryless after a single period}, and the diameter is $1$. Second, the
attractor sizes are exactly powers of two,
\begin{equation}
D^\att_{\max}(N) \;=\; 2^{\lfloor (N+1)/3\rfloor},
\qquad N=6,\ldots,16:\quad 4,4,8,8,8,16,16,16,32,32,32 ,
\end{equation}
which gives $b_\att=2^{1/3}$ \emph{exactly} rather than as R9's fitted value.
The attractor counts run $11,15,20,27,36,48,64,85,113,150,199$ for $W_{28}$ and
$5,7,9,12,16,21,28,37,49,65,86$ for $W_{199}$; both ratios settle at $\rho^2$
per two sites.

\begin{center}
\small
\iftable{tab_r17_shapes_obc0.tex}
\end{center}

\noindent That table also settles the headline question far more strongly than
the largest attractor alone can. It classifies \emph{every} terminal SCC of all
twelve rules at $N=16$ --- $22\,945$ of them --- and the ``cycles'' column is
$\mathbf{0}$ in every row. Not one attractor with more than a single state, in
any of the four classes, is a cycle. (Fixed points are broken out separately
because a lone state with a self-loop is trivially both a cycle and a complete
digraph, and counting it as either would be bookkeeping dressed as evidence.)
The complete count also closes the parent argument numerically: $W_{156}$ and
$W_{198}$ have $190$ complete non-trivial sectors at $N=16$, and $W_{28}$ and
$W_{70}$ have exactly $190$ non-trivial attractors, all of them complete.

\subsection{The spectral gap does not close}

The informative dynamical statement is not the diameter --- a dense graph has a
small diameter for cheap reasons --- but the mixing rate. For the
$\varphi$-classes $|\lambda_2|$ of the row-normalised adjacency runs
\[
0.2786,\; 0.2899,\; 0.2949,\; 0.2973,\; 0.2986,\; 0.2993
\]
as the attractor grows through $21, 55, 144, 377, 987, 2584$ states: flat to
three digits over two decades in $n$, and creeping \emph{up} rather than towards
$1$. The mixing time inside the attractor is $O(1)$ while the attractor itself is
$\Theta(\varphi^N)$. For the plastic class the gap is maximal --- $|\lambda_2|=0$
exactly. Class 2 is the one that does not settle to a single number: $W_{156}$
reads $0.5, 0.597, 0.5, 0.5, 0.5$ and $W_{198}$ reads
$0.333, 0.597, 0.333, 0.5, 0.333$ over $N=8\ldots16$, oscillating inside
$[1/3, 3/5]$ with $N \bmod 4$. It is worth noting that the two differ at all:
they are a reflection pair with identical sizes, densities and reciprocities, and
their out- and in-degree distributions are exact mirrors of each other --- but
$P$ is row-normalised, which is not a reflection-covariant operation, so the
mixing rate of the forward walk need not agree. In none of the four classes does
the gap close as the system grows, which is the point.

\subsection{What grows, and what does not}

The density does fall, and slowly: $0.735, 0.661, 0.594, 0.534, 0.479, 0.431$
over $N=8\ldots16$, a factor $0.900$ per two sites, i.e.\ $\sim0.9487^N$. So the
graphs thin asymptotically, while the mean out-degree grows like $1.534^N$ and
the edge count like $2.481^N$. We report those two bases as fitted; a search over
minimal polynomials of degree $\le3$ with small integer coefficients does not
identify either, unlike every base R9 reports for these rules, and we leave that
open rather than assert an algebraic form we cannot exhibit.

\section{There are two structures here, not four}
\label{sec:parents}

Every one of the twelve rules has a coherent parent obtained by replacing each
reset symbol with the identity, and there are only two parents:
\[
W_{73}\to W_{201},\quad W_{109}\to W_{108};\qquad
W_{28},W_{29},W_{157}\to W_{156},\quad W_{70},W_{71},W_{199}\to W_{198}.
\]
Classes 1 and 2 \emph{are} those parents. The content of this section is that
classes 3 and 4 are as well.

\begin{proposition}
For each of the eight dissipative rules and for $N=8,10,12$, every terminal SCC
of the child is a Krylov sector of its coherent parent --- the same set of
states --- and the graph the child induces on it has \emph{the same edge set} as
the parent's. The reset writes nothing anywhere on its own attractors.
\end{proposition}

\begin{center}
\small
\iftable{tab_r17_parents_obc0.tex}
\end{center}

\noindent Set equality alone would be the weak version of this claim: a child
could inherit the states and still have different edges, since a reset that
fires merges branches. It does not. So the dissipation here creates no new
structure at all; it acts purely as a \textbf{selection rule} on the parent's
sector list, and the whole difference between the four classes is which sectors
survive:

\begin{center}
\begin{tabular}{llll}
\toprule
class & what it keeps & $a_\att$ & $b_\att$ \\
\midrule
$108$, $201$ & every sector & $\rho^2$ & $\varphi$ \\
$73$, $109$ & a $\psi$-sized subfamily, including the largest & $\psi$ &
  $\varphi$ \\
$156$, $198$ & every sector & $\varphi$ & $4^{1/5}$ \\
$28,29,70,71,157,199$ & the sectors that are already complete digraphs & $\rho$ &
  $2^{1/3}$ \\
\bottomrule
\end{tabular}
\end{center}

\noindent The last row is exact for the pure-\rt{D} children: $W_{28}$'s
attractors are \emph{precisely} $W_{156}$'s complete sectors, $36$, $64$ and
$113$ of them at $N=10,12,14$, matching its attractor count series
term for term. The children carrying an \rt{E} keep a strict subset --- $W_{199}$
retains $16$, $28$, $49$ of $W_{198}$'s $36$, $64$, $113$ --- and both still
count at $\rho$ per site. Recording that difference rather than smoothing it is
the honest version: ``the plastic class is the complete-sector subfamily'' is
exactly right for two of the six rules and one containment short for the other
four.

This also explains the coincidence noted in \S1. $b_\att=\varphi$ is shared by
classes 1 and 3 because the largest surviving sector is the same object; the two
classes differ only in $a_\att$, i.e.\ in how many sectors the reset spares, and
$\rho^2$ against $\psi$ is a counting difference within one family of graphs
rather than two different kinds of dynamics.

\section{Which graph measures earn their place}

The question included whether the objects are still too small for
graph-theoretic measures. They are not: $n$ reaches $2584$ with $2.9$M edges,
which is ample for degree distributions, distances, connectivity and spectra.
But two standard measures are uninformative here, and the reason is the density
rather than the size. We would rather say so than present a saturated number as
a result.

\paragraph{Clustering says nothing.} At density $p$, a random digraph already has
transitivity $\approx p$. $W_{201}$ at $N=16$ measures $0.805$ against a
density-matched null of $0.723$: a real but small excess on a scale that is
nearly full. The module therefore never reports transitivity without its null
beside it, and a regression test enforces that.

\paragraph{Connectivity says nothing either.} In every case small enough to
compute it exactly ($n\le200$, since it is $n$ max-flow solves and costs two
minutes already at $n=377$), the vertex and edge connectivities are equal to
each other \emph{and} to the minimum degree: $\kappa_v=\kappa_e=\delta=37$ for
$W_{108}$ at $N=10$, $7$ for $W_{156}$, $7$ for $W_{199}$. That is the generic
value for a dense graph; it describes the sparsest vertex, not the structure.

\paragraph{What does discriminate.} Four things, all of which appear in
Fig.~\ref{fig:topology}: the density and its scaling, the degree profile (a
structured staircase taking $O(\log n)$ distinct values, with the maximum
out-degree equal to $n$ itself --- some states reach the entire attractor in one
period), the reciprocity, which falls steadily from $0.605$ at $n=21$ to $0.319$
at $n=2584$ and so certifies that these graphs are genuinely directed rather
than symmetrised walks, and the spectral gap.

\section{Answering the question as asked}

\begin{enumerate}
\item \textbf{Are they simply cycles?} No, and not marginally. Of the $22\,945$
terminal SCCs the twelve rules have between them at $N=16$, \emph{none} with more
than one state is a cycle. Every one carries a self-loop on every state, has
density of order $1/2$ or exactly $1$, and has diameter $2$ or $1$ (once $3$). A
cycle has density $1/n$ and diameter $n-1$.

\item \textbf{Are they a more connected and extended structure?} More connected,
yes --- dramatically. \emph{Extended}, no: that is the surprise. The number of
states grows exponentially and the graph diameter stays at $2$. These are
compact, densely wired objects, and the monitored walk inside one mixes in $O(1)$
periods no matter how large it is.

\item \textbf{Distinct from the V-free circuits?} Distinct in kind, not in
degree. A V-free rule cannot produce anything but a cycle, because its map is a
function; the branching that makes these graphs two-dimensional is exactly the
Hadamard. And the V-free circuits are not even long: $75$ of $81$ rules have only
fixed points.

\item \textbf{Are they too small for graph measures?} No --- but clustering and
connectivity are uninformative on graphs this dense, and we report both against
their null values rather than as findings. Density scaling, degree profile,
reciprocity and the spectral gap are the measures that separate these graphs
from a random one.

\item \textbf{Unlooked-for, and the main result.} The four classes are two
coherent structures. All eight dissipative rules inherit their parent's sectors
whole, edges included, and differ only in which sectors the reset spares ---
which for the plastic class is exactly the complete ones, and which turns
$b_\att=2^{1/3}$ from a fitted constant into $2^{\lfloor(N+1)/3\rfloor}$.
\end{enumerate}

\section{Reproduction}

\begin{verbatim}
python -m qca_fragmentation.scaling.attractor_topology --bc obc0 --rebuild
python -m pytest qca_fragmentation/tests/test_attractor_topology.py -q
\end{verbatim}

\noindent Everything in this report comes from
\fn{analytics/attractor\_topology\_obc0.json}. Two implementation notes, because
both changed what was feasible to measure. All-pairs distances are computed by
boolean matrix powers rather than by BFS from every node: networkx's $O(nm)$ pass
is $\sim10^{10}$ Python steps at $n=2584$, $m=2.9$M, whereas the diameter being
$2$ means the matrix loop runs twice, each step one BLAS product on a $26$\,MB
\fn{float32} array --- a trade that is only right \emph{because} the answer is
small, and a long-diameter graph would want BFS. Transitivity is likewise taken
as $\operatorname{tr}(B^3)/\sum_i d_i(d_i-1)$ rather than by enumerating
triangles. Connectivity has no such shortcut and is the one measure that is
capped by size rather than computed everywhere.

\end{document}
